\section{The KFD-4/5/6 Discovery Mechanism}
\label{sec:mechanism}

The mechanism is a sequence of gates, not one model call. Each stage preserves
a different authority boundary.

\subsection{KFD-4: replay changes the observation horizon}

For admitted facts $F$, a perspective $P$ defines a projection
$\pi_P(F)$. A transformation $T_{PQ}$ reconstructs a new declared view without
claiming to become the source observer:

\begin{equation}
T_{PQ}: (P, \pi_P(\gamma)) \rightarrow
(Q, \pi_Q(\gamma), J, L),
\end{equation}

where $\gamma$ is a causal path, $J$ is the set of preserved invariants, and
$L$ is declared loss. Source identity, Fact and evidence cuts, known causal
relations, receipts, and authority boundaries may belong to $J$. Relevance,
ordering, natural object boundaries, and available action may change.

Contrastive replay keeps at least two source views distinct. Its purpose is not
to vote on one universal view. It asks which burden becomes natural for the
participant who bears the next consequence. A user may expose an object hidden
from an implementation view; a successor agent may expose a settlement object
hidden from the coordinator who assembled its components.

KFD-4 is one genesis method. Direct situated judgment, anomaly search,
reconstruction analysis, causal-variable discovery, and structural compression
remain first class. This method plurality prevents a successful perspective
story from becoming a ranking theorem.

\subsection{KFD-5: generation is not qualification}

Genesis binds the perspective, current ontology, evidence cut, method,
observation, candidate, and claim boundary. Qualification then asks whether the
candidate deserves independent identity, boundary, authority, lifecycle, or
operations.

For candidate $x$, define $\operatorname{remove}_x(h)$ as a history that erases
the semantic distinction $x$ while preserving other observable facts. The
candidate is non-redundant for decision $d$ when two valid histories exist such
that:

\begin{equation}
\operatorname{remove}_x(h_1)=\operatorname{remove}_x(h_2)
\quad \land \quad
D_d(h_1) \ne D_d(h_2).
\end{equation}

This deletion witness does not privilege a name or physical representation. A
rival model passes if it preserves equivalent information with lower total
burden. KFD-5 therefore requires alternatives, minimum closure, deletion and
fuse tests, falsifiers, dogfood under load, residual risks, and one of five
outcomes: accept, provisional, reject, subsume, or no new Primitive justified.

\subsection{KFD-6: a bounded experimental loop}

At loop step $n$, let $E_n$ be an immutable causal-experience cut, $O_n$ the
declared ontology, $G_n$ plural generation experiments, $C_n$ a shared-budget
method comparison, $Q_n$ KFD-5 qualification, $H_n$ held-out or independent
evaluation, and $A_n$ promotion authority. Draft KFD-6 proposes:

\begin{equation}
(E_n,O_n) \rightarrow G_n \rightarrow C_n \rightarrow Q_n
\rightarrow H_n \rightarrow A_n \rightarrow O_n \;\text{or}\; O_{n+1}.
\end{equation}

The conservative baselines are part of the mechanism:

\begin{itemize}[leftmargin=*]
  \item \textbf{fixed ontology:} improve prediction, retrieval, or workflow
        while forbidding new first-class objects;
  \item \textbf{no new Primitive:} permit diagnosis and local abstraction but
        require the final decision to retain $O_n$; and
  \item \textbf{flattened history:} use anonymous events or endpoint facts
        under the same resource budget.
\end{itemize}

Candidates are evaluated for yield, false-candidate rate, ontology distance,
compression, intervention value, transfer, falsifiability, cost, and
auditability. Negative results remain part of the cut. Generated examples
cannot be the only evidence. The generator cannot be the sole verifier or
promotion authority.

\subsection{The authority flow}

The complete mechanism assigns a narrow authority to each layer:

\begin{center}
\begin{tabular}{p{0.20\linewidth}p{0.68\linewidth}}
\toprule
Layer & Authority \\
\midrule
Episode corpus & Evidence carrier with declared capture and stewardship; not
an ontology decision. \\
KFD-4 replay & Declared projection and comparison; not source mutation or
candidate certification. \\
Genesis method & Proposes a candidate and disconfirming test; not truth or
promotion. \\
KFD-5 qualification & Judges bounded load-bearing utility against facts and
alternatives; not universal adoption. \\
Held-out evaluator & Tests transfer and false candidates outside the discovery
cut; not value or legal authority. \\
Promotion authority & Admits or rejects ontology change inside a declared
scope; does not rewrite prior Episodes. \\
\bottomrule
\end{tabular}
\end{center}

This separation is the central anti-self-certification property. It lets an
agent contribute ontology change without turning model confidence into social
or operational authority.
